%========================================================
%  Appendix A. Band-limited kernels and BS sandwiches
%========================================================
\appendix
\section{Band-limited kernels and Beurling–Selberg sandwiches}\label{sec:appendix-kernels}

Throughout $A\ge 10$, $H=(\log X)^A$, and implicit constants may depend on $(\varepsilon,A,C_W)$ only.

%----------------------------------------

\subsection*{Kernel elaboration}
\begin{lemma}[Admissible band-limited kernel]\label{lem:kernel}
For each $H\ge2$ define
\[
\widehat K_H(\xi)\ :=\ A_H\,|\xi|\,\mathbf{1}_{\{|\xi|\le 1/H\}}(\xi),
\qquad A_H:=H.
\]
Let
\[
K_H(t)\ :=\ \int_{-1/H}^{1/H}\widehat K_H(\xi)\,e\!\left(\frac{\xi t}{X}\right)\,d\xi,
\qquad e(z):=e^{2\pi i z}.
\]
Then $\widehat K_H(\xi)\ge0$, is even and compactly supported in $[-1/H,1/H]$, and
\begin{equation}\label{eq:KH-moments-lemma}
\int_{-1/H}^{1/H}\!\widehat K_H(\xi)\,d\xi\ \ll\ H,\qquad
\int_{-1/H}^{1/H}\!\frac{\widehat K_H(\xi)}{|\xi|}\,d\xi\ \ll\ H\log H,
\end{equation}
with absolute implied constants. Moreover $K_H$ is real, even, and nonnegative-definite.
\end{lemma}

\begin{proof}
The nonnegativity and support are immediate from the definition; evenness follows from $|\xi|$.
Since $\widehat K_H$ is integrable and real-even, $K_H$ is real-even, and Bochner's theorem yields nonnegative-definiteness.

For the moment bounds, compute explicitly:
\[
\int_{-1/H}^{1/H}\widehat K_H(\xi)\,d\xi
= 2\int_{0}^{1/H} A_H\,\xi\,d\xi
= \frac{A_H}{H^2}
= \frac{1}{H}
\ \ll\ H.
\]
Likewise,
\[
\int_{-1/H}^{1/H}\frac{\widehat K_H(\xi)}{|\xi|}\,d\xi
= 2\int_{0}^{1/H} A_H\,d\xi
= \frac{2A_H}{H}
= 2
\ \ll\ H\log H.
\]
This proves \eqref{eq:KH-moments-lemma}.
\end{proof}

\noindent\emph{Kernel moments.}
All integrals in $\xi$ use the admissible kernel moments \eqref{eq:K-moments}, proved in Lemma~A.1 ((A.2)–(A.3)) of Appendix~A for our explicit model $\widehat K_H(\xi)=H\,\phi(H\xi)$.

\subsection*{A concrete positive band-limited majorant}\label{subsec:KH}

Fix an even $\eta\in C_c^\infty([-1,1])$ with $0\le \eta\le 1$, $\eta(0)=1$. Define
\[
  \widehat{K_H}(\xi)\ :=\ H\,\eta(\xi H)^2 \qquad (\xi\in\R),\qquad
  K_H(t)\ :=\ \int_{\R}\widehat{K_H}(\xi)\,e(-\xi t)\,d\xi.
\]
Then $K_H$ is real, even, nonnegative definite, band-limited to $|\xi|\le 1/H$, and satisfies:
\begin{align}
  \int_{\R} K_H(t)\,dt \,=\, \widehat{K_H}(0)\,=\,H, \qquad
  \|\widehat{K_H}\|_{L^1(\R)} \,\ll\, 1 \,\le\, H, \label{eq:KHmassL1}\\
  \forall\,\tau\in(0,1/H]\!:\ \ 
  \int_{\tau\le |\xi|\le 1/H}\frac{|\widehat{K_H}(\xi)|}{|\xi|}\,d\xi \ \ll\ H\,\log\frac{1}{\tau H}.
  \label{eq:KHlogmom}
\end{align}
\emph{Remarks.} (i) \eqref{eq:KHmassL1} gives the $L^1$ size in frequency used in TT$^\ast$.  
(ii) Taking $\tau=H^{-2}$ yields the logarithmic moment bound $ \ll H\log H$ used in §SSU.  
(iii) If one prefers an explicit closed form, take $\eta=\vartheta\!\ast\!\vartheta$ with $\vartheta$ a $C^\infty$ bump on $[-\tfrac12,\tfrac12]$; then $\widehat{K_H}=H(\eta(\xi H))^2$ remains $\ge 0$ and compactly supported.

%----------------------------------------
\subsection*{Beurling–Selberg majorant/minorant for \texorpdfstring{$[-H,H]$}{[-H,H]}}\label{subsec:BS}

There exist even, real $M_H,m_H\in L^1(\R)$ with $\supp\widehat{M_H},\supp\widehat{m_H}\subset[-1/H,1/H]$ such that
\begin{equation}\label{eq:BSsandwich}
  m_H(t)\ \le\ \mathbf{1}_{[-H,H]}(t)\ \le\ M_H(t)\qquad (t\in\R),
\end{equation}
and the following hold:
\begin{align}
  \int_{\R}\!\big(M_H(t)-m_H(t)\big)\,dt \ \ll\ 1, \qquad
  \|\widehat{M_H}\|_{1}+\|\widehat{m_H}\|_{1}\ \ll\ H, \label{eq:BSl1}\\
  \int_{|\xi|\le 1/H}\!\frac{|\widehat{M_H}(\xi)|+|\widehat{m_H}(\xi)|}{|\xi|}\,d\xi \ \ll\ H\log H. \label{eq:BSlog}
\end{align}
\emph{Remarks.} This is the classical Beurling–Selberg extremal pair for an interval, scaled to bandwidth $1/H$. One may cite, e.g., Vaaler’s explicit construction or standard Selberg–Beurling references. The $L^1$ and logarithmic moment bounds \eqref{eq:BSl1}–\eqref{eq:BSlog} are what the arguments in §§PSB/BG/AO use.

\begin{proof}
Fix $(d,n)$. The condition $(d,n)\in T(a/q,\Delta)$ is $|d'n-dn'|\le H$ with $d'/d$ in a $\asymp Q^{-2}$ window around $a/q$. Two distinct slopes $a/q\ne a'/q'$ give disjoint windows since $|a/q-a'/q'|\gg Q^{-2}$ and the short-axis thickness in slope is $\ll Q^{-2}$. Hence $(d,n)$ belongs to $O(1)$ tubes. The pairwise intersections for two slope families lie in a rectangle of long side $\asymp X$ and short side $\asymp X/H$, hence $O(1)$ per short-axis run. Summing the local bounds and invoking Proposition~\ref{prop:incidence} for the global counting yields the stated inequalities.
\end{proof}

%----------------------------------------
\subsection*{Failed hypothesis: A fixed nonnegative banked pin}\label{subsec:fixed-pin}

Initially, our Route B involved a fixed nonnegative banked pin instead of smoothed bank projection, as follows. It did not work.

Let $\mathfrak A=\bigcup_{j=1}^{J}\mathfrak a_j$ be the AO$^+$ bank: disjoint arcs of width $\asymp 1/H$, $J\asymp Q^2$.
Shrink each to its middle third $\mathfrak a'_j\subset\mathfrak a_j$ so that the Minkowski sum still lies in $\mathfrak A$.
Choose $\phi_j\in C_c^\infty(\mathfrak a'_j)$ with $0\le\phi_j\le 1$ and $\int \phi_j\asymp 1/H$, and set
\[
  \Phi(\xi):=\sum_{j=1}^J \phi_j(\xi),\qquad \Psi(\xi):=\sqrt{\Phi(\xi)},\qquad
  k:=\mathcal F^{-1}\Psi,\qquad K:=|k|^2.
\]
Then $K(n)\ge 0$, $\sum_{n\in\Z} K(n)=\big(\int\Psi\big)^2\asymp 1$, and $\widehat K=\Psi\ast\Psi$ is supported in $\mathfrak A$ (by the middle–thirds choice). Moreover,
\begin{equation}\label{eq:pin-boost}
  \forall\,\varepsilon\in(0,1):\quad
  K_{\mathrm{pin}}\ :=\ \underbrace{K*K*\cdots*K}_{r\ \text{times}}\ \ \text{with } r\asymp \log(1/\varepsilon)
  \ \Rightarrow\ 
  K_{\mathrm{pin}}(0)\ge 1-\varepsilon,\ \ \sum_{n\neq 0}K_{\mathrm{pin}}(n)\le \varepsilon,
\end{equation}
and $\supp\,\widehat{K_{\mathrm{pin}}}\subset \mathfrak A$.
This is the “fixed positive pin’’ formerly used in §\ref{sec:center-uniform}; \eqref{eq:pin-boost} follows from standard large–deviation for symmetric convolutions (or Young’s inequality).

The pin was abandoned and replaced with bank projection when the pin was shown to be inadequate to prove full Goldbach. A two-pin solution may be necessary if we were to attempt to solve the twin primes conjecture, but that goes beyond the scope of this paper.

%----------------------------------------
\subsection*{Bernstein and sampling on a \texorpdfstring{$\ll H$}{<<H} grid}\label{subsec:bernstein-LS}

\begin{lemma}[Bernstein]\label{lem:Bernstein}
If $f\in L^\infty(\R)$ has $\supp\widehat f\subset[-\Lambda,\Lambda]$, then $\|f'\|_\infty\le 2\pi \Lambda\,\|f\|_\infty$. In particular,
if $\supp\widehat f\subset[-1/H,1/H]$ and $|x-y|\le \Delta$, then
$|f(x)-f(y)|\le (2\pi\,\Delta/H)\,\|f\|_\infty$.
\end{lemma}

\begin{lemma}[Band-limited sampling on a thick grid]\label{lem:LS-grid}
Let $f$ be band-limited to a measurable $\Sigma\subset[-1/H,1/H]$ with $|\Sigma|\ll Q^2/H$. 
Let $\mathcal G=\{X+j\Delta\}\cap[X,2X]$ with $\Delta=H/C_0$ and $C_0$ large.
Then 
\[
  \|f\|_{L^\infty([X,2X])}\ \ll\ (\log X)^C\ \sup_{y\in\mathcal G}|f(y)|.
\]
If, in addition, $f(y)\ge \alpha$ on at least $(1-\delta)$ of the grid points with $\delta\le \log^{-B}X$, then
$\inf_{[X,2X]} f \ge \tfrac12 \alpha$ for $B$ large (constants depend only on $(C_0,A)$).
\end{lemma}

\emph{Sketch.} Quantitative Logvinenko–Sereda (Kovrijkine/Nazarov) for band-limited functions on thick sets, with the bandwidth and $|\Sigma|$ dependence absorbed into $(\log X)^C$; then apply Lemma~\ref{lem:Bernstein} to pass from grid to continuum.

%----------------------------------------
\subsection*{Beurling–Selberg sandwich inequality}\label{subsec:BS-ineq}

For any arithmetic function $F(n)$ supported on $[X,2X]$ and any $y\in[X,2X]$,
\[
  \sum_{|n-y|\le H} F(n)\ \ge\ \sum_{n} F(n)\,m_H(n-y),\qquad
  \sum_{|n-y|\le H} F(n)\ \le\ \sum_{n} F(n)\,M_H(n-y),
\]
with $m_H,M_H$ as in \S\ref{subsec:BS}. The error between the two sides integrates to $O(1)$ uniformly in $y$ by \eqref{eq:BSl1}.
This is the version used in §§no-long-gap and center–uniform.

\bigskip

\noindent\textbf{References for Appendix A.} 
Classical constructions of Beurling–Selberg extremals and Vaaler’s explicit majorants/minorants for intervals; standard Bernstein inequality; quantitative Logvinenko–Sereda (Nazarov; Kovrijkine).
